\section{Introduction}
\label{sec:rewrite-intro}

Symmetry block diagonalizes a Hamiltonian, but block diagonalization alone
does not determine which eigenstates a protocol can occupy.  If an initial
state is fixed by every preserved symmetry, evolution is confined to the
joint fixed space.  Even that statement can be loose: the algebra generated
by the endpoint Hamiltonians may act cyclically only on a proper subspace of
the joint fixed space.  This cyclic subspace is a common invariant envelope,
not a claim that every vector in it is reachable under one prescribed
schedule.  A small global gap can then describe a crossing that the protocol
neither couples to nor detects.

Set Cover, whose minimum-cardinality version is NP-hard
\cite{karp1972reducibility}, admits the usual one-binary-variable-per-base Ising formulation,
and it also appears among the constrained examples motivating quantum
alternating-operator ansatzes~\cite{Lucas2014,Hadfield2019}.  Separately,
qubit-efficient variational encodings show that logarithmic or otherwise
compressed register counts are not by themselves a novelty
claim~\cite{TanEncoding2021,GlosKrawiecZimboras2022,Chancellor2019}.  Our
multi-register address encoding belongs to this broad line; no result below
depends on claiming a new qubit count.

Classical objective symmetries can confine quantum approximate optimization
algorithm (QAOA) dynamics and impose equal measurement probabilities on
symmetry-related strings
\cite{ShaydulinHadfieldHoggSafro2021}.  Quantum walks on automorphism
quotients likewise identify a smaller invariant evolution for suitable
initial states~\cite{KroviBrun2007}, while Krylov methods use repeated
generator action to identify a minimal dynamical subspace
\cite{Saad1992}.  In adiabatic computation, symmetry-adapted path design was
studied by Schaller and Sch\"utzhold~\cite{SchallerSchuetzhold2010}, and
symmetry-restricted effective gaps of stoquastic Hamiltonians were connected
to classical simulability by Bringewatt and Jarret
\cite{BringewattJarret2020}.  These results make an effective sector
physically relevant but do not identify the joint fixed space with the
smaller endpoint-algebra envelope considered here.  Conversely, modern
time-dependent Krylov constructions adapt the basis directly to a driven
trajectory~\cite{TakahashiDelCampo2025}.  We use a different, exact object:
the common invariant envelope of the endpoint family.  A specified schedule
traces a potentially smaller subset of it, so no controllability converse is
assumed.

Finally, symmetric discriminants of reversible Markov chains and annealing
paths with Gibbs-amplitude ground states are established tools in quantum
walk and quantum simulated-annealing constructions
\cite{Szegedy2004,SommaBoixoBarnumKnill2008}.  Quantitative adiabatic
computation originates with the interpolation framework of
Farhi~et~al.~\cite{Farhi2000}; modern reviews emphasize that runtime claims
require a specified path and access model~\cite{AlbashLidar2018}.
Quantitative statements require an isolated spectral band together with
derivative and gap control~\cite{JansenRuskaiSeiler2007}.  We use the discriminant only to
construct a rigorously gapped witness path and make no speedup claim.

Minimum Set Cover (MSC) provides a concrete setting in which these three
spaces differ.  With $n_B$ candidate bases, a $k$-register encoding
represents an ordered candidate by $k\lceil\log_2n_B\rceil$ address qubits.
Permuting the registers gives an action of the symmetric group $S_k$, while
colour-preserving automorphisms of the incidence relation permute the base
labels~\cite{McKayPiperno2014}.  Register symmetry alone and a
final-Hamiltonian sector gap are insufficient for an adiabatic or otherwise
path-dependent statement: one must use the faithful joint action, the common
cyclic envelope selected by the initial state, and the minimum isolation gap
of the band tracked along an explicit path.

This work makes that distinction its central question.  Its three main
contributions are:
\begin{enumerate}
\item an orbit description of the $k$-register encoding under
$S_k\times G_B$, where $G_B$ is the faithfully represented group of base
automorphisms, $\mathcal H_{\rm sym}$ is the subspace fixed by this joint
action, and $\mathcal K_u$ is the common invariant cyclic envelope generated
from the uniform initial state by the endpoint algebra; this gives invariant covering
penalties, exact P\'olya--Redfield orbit counts
\cite{Redfield1927,Polya1937}, and a constructive distinction between these
two spaces;
\item a sector-localization proposition with its kinetic floor, together
with a general orbit-resolved hopping-cancellation theorem and a separate
stability theorem for transverse inter-sector crossings; and
\item an even-cycle illustration of the cancellation theorem and,
separately, a parent protocol built from Metropolis rates on Johnson-graph
moves~\cite{Metropolis1953,BrouwerCohenNeumaier1989}, whose cyclic-envelope gap
is uniformly polynomial, yielding a conditional polynomial-time adiabatic
corollary.
\end{enumerate}

The last result concerns a protocol-dependent invariant gap, not computational
advantage.  The controlled cycle family is classically easy, and the proven
exponent is conservative.  Its role is to separate the gap question cleanly
into symmetry, geometry, and kinetics, without inferring asymptotic behavior
from finite-size data.

Neither orbit reduction, P\'olya enumeration, compact addressing, nor a
Markov discriminant is claimed as new in isolation.  The contribution is
the protocol-dependent link between orbit potentials, group-fixed spectra,
cyclic-envelope closure, tracked bands, sector localization, and dark
multiplicity, together with a separately controlled parent path.

We use the following notation throughout.  For an operator $X$,
$\Ran X$ and $\spec X$ denote its range and spectrum.  The symbol
$\|\cdot\|$ denotes the Hilbert-space norm on vectors and the induced
operator norm on operators.  The gap of a finite-dimensional Hermitian
operator is the distance from its lowest eigenspace to the next distinct
eigenvalue; $\lambda_{\min}(X|_{\mathcal L})$ is the lowest eigenvalue of
the compression of $X$ to $\mathcal L$.  For positive $g$, the notation
$f=O(g)$ means $|f|\leq Cg$ for a constant $C>0$ independent of the
asymptotic variable in the stated limit.
